\section{Architecture}
\label{sec:architecture}

\subsection{The projection engine}
The unit of physical work is one $1280\times320$ projection: the decoder feed-forward output matrix, the largest shape in the model. \begin{figure*}[t]
\centering
\resizebox{\textwidth}{!}{%
\begin{tikzpicture}[font=\footnotesize,>=stealth,node distance=6mm,
  box/.style={draw,rounded corners=1.5pt,minimum height=9mm,minimum width=17mm,align=center,inner sep=3pt},
  new/.style={box,fill=blue!6}]
\node[box] (in) {activation $a$\\(signed 9-bit)};
\node[box,right=of in] (pp) {16 partials $a\!\cdot\!k$\\$\times$10 local copies};
\node[box,right=9mm of pp] (lanes) {320 lanes\\select $+$ add\\27-bit accumulators};
\node[new,right=9mm of lanes] (grp) {10 group\\registers\\(look-ahead)};
\node[new,right=of grp] (corr) {$-\,\Sigma a\cdot z$\\(registered)};
\node[new,right=of corr] (pipe) {elastic\\register};
\node[box,right=of pipe] (fp) {FP32 affine\\HardFloat,\\3 stages};
\node[box,below=8mm of pp] (arr) {index array $1280\times320\times4$\,b\\(external, mask-fixed)};
\node[box,right=9mm of arr] (book) {16-entry books\\$\to$ UINT8};
\draw[->] (in)--(pp); \draw[->] (pp)--(lanes);
\draw[->] (arr)--(book); \draw[->] (book.north) -- (book.north |- lanes.south);
\draw[->] (lanes)--(grp); \draw[->] (grp)--(corr); \draw[->] (corr)--(pipe); \draw[->] (pipe)--(fp);
\draw[->] (fp.east) -- ++(6mm,0) node[right,align=left]{FP32\\out};
\node[below=2mm of corr,font=\scriptsize,text=blue!60!black] {shaded: added for slow-corner closure (\S\ref{sec:closure})};
\end{tikzpicture}}
\caption{The $1280\times320$ projection engine. The index array sits outside the engine's physical top. One projection takes 4,165 cycles: 1,280 rows of LOAD/INDEX/ACCUM, one PREP cycle, 320 output cycles and four pipeline-fill cycles.}
\label{fig:engine}
\end{figure*}

The engine (\cref{fig:engine}) is a five-state machine:
\begin{description}
\item[LOAD] accepts one signed 9-bit activation, adds it to a running activation sum, and registers its 16 partial products.
\item[INDEX] reads one 1,280-bit row of the index array (320 four-bit selections) and maps each index through its book to a UINT8 representative.
\item[ACCUM] adds, in each of 320 lanes, the two selected partial products (one shifted) into a signed 27-bit accumulator.
\item[PREP] (added in \cref{sec:closure}) fills per-group output registers before output begins.
\item[OUTPUT] emits one column per cycle. Each value is the accumulator minus $\text{activation sum}\times\text{zero point}$, passed through a three-stage FP32 multiply--add (Berkeley HardFloat~\cite{hauser_hardfloat}) that applies the activation scale, the weight scale and the bias with independent round-to-nearest-even steps.
\end{description}
One projection takes $3\times1280$ cycles of LOAD/INDEX/ACCUM and 320 output cycles, plus pipeline fill: 4,165 cycles in the final design (4,163 originally). That is 409,600 multiply-accumulates, an average of 98 MAC per cycle.

Two structural changes made the engine physically implementable:
\begin{enumerate}
\item \textbf{Fanout.} Row-index and partial-product nets originally drove up to \num{11359} loads. We duplicated the row registers ten times and capped synthesis fanout at 16.
\item \textbf{Local partials.} Each 32-lane group gets its own copy of the partial-product registers, so the 208-bit partial-product bus no longer crosses the die (\cref{sec:closure}).
\end{enumerate}
Both changes are cycle-identical on the full test bench.

\subsection{From one engine to a recognizer}
The full recognizer RTL connects:
\begin{itemize}
\item the PCM input and frontend;
\item the six-layer encoder;
\item an adapter that produces cross-attention keys and values;
\item a greedy six-layer decoder that ends at the end-of-sequence token.
\end{itemize}
Everything runs on one common clock. Each weight service answers lookups through index and book ROMs. The top level has no weight inputs: its 31 runtime inputs are audio, clocks, resets, handshakes and configuration.

\paragraph{Throughput.}
A dense one-pass estimate of the workload is about 0.6\,GMAC/s:
\begin{itemize}
\item the encoder: about 7.4\,M MAC per frame at 50 frames/s;
\item the decoder: about 21.5\,M MAC per token at about 4 tokens/s;
\item the keys and values produced for cross-attention.
\end{itemize}
One engine at \SI{150}{ns} sustains 0.66\,GMAC/s. However, the native streaming controller re-processes encoder frames and regenerates cross-attention keys and values. Measured on real traces, that controller needs 10--90 engine-equivalents at a \SI{120}{ns} clock. Real-time operation therefore depends as much on scheduling as on raw MAC rate, and we treat it as open (\cref{sec:limitations}).

\subsection{Chip plan}
All 98 matrices map to 81 pitch-matched index banks (\SI{174.7}{mm^2}) plus four bank-driver control strips (\SI{56.6}{mm^2}). Abstract floorplans with power grids give two candidate chips:
\begin{itemize}
\item \textbf{One die} of $25\times31$\,mm, which needs a denser SRAM than the open GF180 macro (we assumed \SI{12}{\micro m^2}/bit);
\item \textbf{Two dies} of $25\times29$\,mm, built only from existing open parts, with no change to the numerics.
\end{itemize}
The two-die variant routes and passes IR drop at the measured engine power, but fails an earlier, more pessimistic assumed-load IR screen.

About 1.09\,MB of mutable state (activation banks, key/value caches, history and scratch) dominates the remaining area. Synthesized as flip-flops, it alone puts the logic at \SI{1198}{mm^2} (\cref{sec:results}). Mapping it onto SRAM macros is the largest open area item.
